\setchapterpreamble[u]{\margintoc}
\chapter{Limit Theorems}
\labch{limit-theorems}

\sources{Bertsekas and Tsitsiklis, \emph{Introduction to Probability}, 2nd ed.\
\cite{bertsekas2008probability}, Chapter~5; MIT 6.041 \cite{ocw6041},
Lectures~19--20.}

Two theorems justify the entire practice of learning from a sample: averages
converge to what they estimate, and the error of the average has a predictable
shape and size.

\section{Two inequalities}
\labsec{inequalities}

\begin{proposition}[Markov]
If \(X\ge0\) and \(a>0\) then \(\Prob(X\ge a)\le\E[X]/a\).
\end{proposition}

\begin{proof}
\(\E[X]\ge\E[X\cdot\mathbf{1}\{X\ge a\}]\ge a\,\Prob(X\ge a)\), using
non-negativity in the first step.
\end{proof}

\begin{proposition}[Chebyshev]
If \(\E[X]=\mu\) and \(\Var(X)=\sigma^{2}\) then for \(a>0\),
\[
	\Prob\bigl(|X-\mu|\ge a\bigr)\le\frac{\sigma^{2}}{a^{2}}.
\]
\end{proposition}

\begin{proof}
Apply Markov to the non-negative variable \((X-\mu)^{2}\) with threshold
\(a^{2}\).
\end{proof}

\marginnote{These bounds are crude --- they use only two moments --- but they
hold for every distribution, which is exactly what is wanted when the
distribution is unknown.}

\section{The law of large numbers}
\labsec{lln}

\begin{theorem}[Weak law]
Let \(X_1,X_2,\ldots\) be i.i.d.\ with mean \(\mu\) and finite variance
\(\sigma^{2}\). Then for every \(\varepsilon>0\),
\[
	\Prob\bigl(|\bar X_n-\mu|\ge\varepsilon\bigr)\longrightarrow0
	\qquad\text{as }n\to\infty .
\]
\end{theorem}

\begin{proof}
By \refsec{variance}, \(\Var(\bar X_n)=\sigma^{2}/n\). Chebyshev gives
\[
	\Prob\bigl(|\bar X_n-\mu|\ge\varepsilon\bigr)\le\frac{\sigma^{2}}{n\varepsilon^{2}},
\]
which tends to zero.
\end{proof}

\begin{remark}
Training loss is an average over the sample, and test loss is an average over a
different sample from the same distribution. The law is what allows either to
stand in for the expected loss --- and its reliance on the i.i.d.\ hypothesis is
what makes distribution shift so damaging.
\end{remark}

\section{The central limit theorem}
\labsec{clt}

\begin{theorem}
With the same hypotheses and \(\sigma^{2}>0\),
\[
	\frac{\bar X_n-\mu}{\sigma/\sqrt n}
	\ \xrightarrow{\ d\ }\ \text{Gaussian}(0,1).
\]
\end{theorem}

The law of large numbers says the error vanishes; the central limit theorem says
it vanishes at rate \(1/\sqrt n\) and that the rescaled error is Gaussian
regardless of the distribution the samples came from.

\marginnote{This is the honest reason Gaussian noise is assumed so often. It is
not that measurements are Gaussian by nature, but that an accumulation of many
small independent contributions is approximately so.}

\begin{example}
Estimating an accuracy \(p\) from \(n\) independent test examples gives
\(\Var(\bar X_n)=p(1-p)/n\). With \(p\approx0.9\) and \(n=1000\), the standard
deviation is about \(0.0095\), so two systems differing by half a percentage
point on such a test set are not distinguishable by that test set alone.
\end{example}

\section{What the theorems do not say}
\labsec{limit-caveats}

Both require independence and identical distribution, and the central limit
theorem additionally requires finite variance. Heavy-tailed quantities can
violate the last condition, and time-ordered data routinely violate the first.
Where the hypotheses fail, the conclusions are not merely weaker --- they can be
false, and the confidence intervals built on them are then decorative.
